\documentclass[11pt,a4paper]{article}

% ---------------------------------------------------------------------------
% Paquetes
% ---------------------------------------------------------------------------
\usepackage[utf8]{inputenc}
\usepackage[T1]{fontenc}
\usepackage{lmodern}
\usepackage[spanish]{babel}
\usepackage[margin=2.5cm]{geometry}
\usepackage{amsmath}
\usepackage{graphicx}
\usepackage{float}
\usepackage{booktabs}
\usepackage{array}
\usepackage{xcolor}
\usepackage{listings}
\usepackage{hyperref}

\hypersetup{colorlinks=true, linkcolor=blue!50!black, urlcolor=blue!50!black}

% ---------------------------------------------------------------------------
% Atajos
% ---------------------------------------------------------------------------

% Hueco para completar a mano con el dato medido.
\newcommand{\dato}[1][2cm]{\underline{\hspace{#1}}}

% Microsegundos sin depender de paquetes de unidades.
\newcommand{\us}{\ensuremath{\mu}s}

% Placeholder de captura: el documento COMPILA sin las imagenes. Cuando
% tengas la foto del osciloscopio, reemplazá el bloque por el \includegraphics
% que figura adentro del recuadro.
\newcommand{\captura}[1]{%
  \begin{figure}[H]
  \centering
  \fbox{\parbox[c][4.5cm][c]{0.82\textwidth}{\centering
    \itshape Captura pendiente\\[6pt]
    \normalfont\small #1\\[10pt]
    \footnotesize Reemplazar este recuadro por:\\
    \texttt{\textbackslash includegraphics[width=0.82\textwidth]\{img/archivo.png\}}}}
  \caption{#1}
  \end{figure}}

\lstdefinestyle{cstyle}{
    language=C, basicstyle=\ttfamily\small,
    keywordstyle=\color{blue!60!black}\bfseries,
    commentstyle=\color{green!45!black}\itshape,
    backgroundcolor=\color{gray!8}, frame=single,
    rulecolor=\color{gray!55}, showstringspaces=false, tabsize=2,
    breaklines=true, captionpos=b, xleftmargin=1.2em
}
\lstset{style=cstyle}

\title{\textbf{Medición de Tiempos de Interrupción}\\[4pt]
\large Guía de ensayos y análisis de resultados --- FRDM-K64F}
\author{Santiago Riverós \and Ignacio Sammartino \and Rodrigo Devesa \and Leandro Yiu}
\date{TP1 --- Sistemas Embebidos}

\begin{document}

\maketitle

\begin{abstract}
\noindent
Este documento define la campaña de mediciones con osciloscopio sobre el
firmware del sistema de control de acceso, y contiene el análisis de los
resultados. Cada ensayo indica qué se mide, cómo se configura el equipo y
--- sobre todo --- qué conclusión de diseño permite extraer. El objetivo no
es sólo caracterizar tiempos, sino \emph{verificar experimentalmente} tres
cosas: que el sistema cumple sus restricciones de tiempo real, cuánto cuestan
las abstracciones de la arquitectura en capas, y que la asignación de
prioridades de interrupción es correcta.
\end{abstract}

\tableofcontents
\newpage

% =============================================================================
\section{Método}
% =============================================================================

\subsection{El pin de medición}

El firmware incluye un módulo de instrumentación (\texttt{source/test\_isr\_cfg.h})
que pone el pin \texttt{PIN\_TEST\_ISR} en alto al entrar a una rutina y en
bajo al salir. De ese modo, en el osciloscopio:

\begin{itemize}
  \item el \textbf{ancho del pulso} es el tiempo de ejecución de la rutina;
  \item el \textbf{período entre flancos de subida} es cada cuánto se ejecuta;
  \item el \textbf{duty} es el porcentaje de CPU que consume.
\end{itemize}

Sólo se puede medir una rutina por vez, porque hay un único pin físico. La
selección se hace cambiando una línea y recompilando:

\begin{lstlisting}[caption={Selección de la rutina a medir en \texttt{test\_isr\_cfg.h}.}]
#define TEST_ISR_TARGET             TEST_ISR_DISPLAY_REFRESH
\end{lstlisting}

La instrumentación se resuelve en tiempo de compilación: los puntos de
medición no seleccionados expanden a \texttt{((void)0)} y no generan ni una
instrucción. El binario con \texttt{TEST\_ISR\_NONE} es idéntico en tamaño al
que tiene la medición activa, de modo que medir no altera lo medido más allá
de las dos escrituras al GPIO.

\subsection{Condiciones del ensayo}

\begin{table}[H]
\centering
\begin{tabular}{@{}ll@{}}
\toprule
\textbf{Parámetro} & \textbf{Valor} \\
\midrule
Microcontrolador        & MK64FN1M0VLL12 (Cortex-M4) \\
Frecuencia de core      & 100 MHz (1 ciclo = 10 ns) \\
Optimización            & \texttt{-O0} (build de Debug) \\
Frecuencia del SysTick  & 1 kHz (período 1 ms) \\
Pin de medición         & \texttt{PIN\_TEST\_ISR} = PTC17 \\
Osciloscopio            & \dato[4cm] \\
Sonda / atenuación      & \dato[4cm] \\
\bottomrule
\end{tabular}
\caption{Condiciones bajo las que se realizaron todas las mediciones.}
\end{table}

\textbf{Importante:} los tiempos corresponden a una compilación sin optimizar.
Con \texttt{-O2} los valores absolutos bajarían de forma apreciable; las
relaciones entre ellos, que son la base de las conclusiones, se mantienen.

\subsection{Configuración del osciloscopio}

\begin{enumerate}
  \item Punta en PTC17, masa a cualquier GND de la placa.
  \item Acoplamiento DC, disparo por \textbf{flanco de subida} en $\approx$1,5 V.
  \item Para medir el \textbf{ancho}: base de tiempo corta, de modo que el
        pulso ocupe buena parte de la pantalla. Usar la medición automática
        de \textit{ancho de pulso positivo}.
  \item Para medir el \textbf{período}: base de tiempo larga, hasta ver varios
        pulsos. Usar la medición automática de \textit{período} o
        \textit{frecuencia}.
  \item Anotar siempre las dos cosas: con ancho y período se calcula el duty.
\end{enumerate}

\subsection{Cálculos a aplicar}

Para cada rutina medida, con ancho $t_{on}$ y período $T$:

\begin{center}
\begin{tabular}{@{}lll@{}}
\toprule
\textbf{Magnitud} & \textbf{Fórmula} & \textbf{Qué significa} \\
\midrule
Frecuencia de ejecución & $f = 1/T$              & cada cuánto corre \\
Carga de CPU            & $D = t_{on}/T \cdot 100\%$ & cuánta CPU consume \\
Ciclos de reloj         & $N = t_{on} \cdot 100\,\text{MHz}$ & costo en instrucciones \\
\bottomrule
\end{tabular}
\end{center}

\newpage
% =============================================================================
\section{Ensayos}
% =============================================================================

Los ensayos están ordenados a propósito: M1 establece la base de tiempo, M2--M4
desglosan sus componentes, y M5--M8 atacan el camino crítico del lector de
tarjetas. Conviene hacerlos en ese orden.

% -----------------------------------------------------------------------------
\subsection{M1 --- SysTick completo}
% -----------------------------------------------------------------------------

\textbf{Configuración:} \texttt{TEST\_ISR\_TARGET = TEST\_ISR\_SYSTICK\_TOTAL}

\textbf{Objetivo.} Caracterizar la interrupción periódica más importante del
sistema y verificar la restricción de tiempo real fundamental: la ISR del
SysTick tiene que terminar \emph{siempre} antes del próximo tick. Si no, el
sistema pierde ticks y todos los tiempos del firmware (antirrebote del
encoder, temporizadores de la FSM, refresco del display) se degradan.

\textbf{Procedimiento.} Con el sistema en reposo (animación de StandBy
corriendo), medir ancho y período. Verificar que el período sea 1,000 ms.

\captura{M1: pulso de la ISR completa del SysTick.}

\begin{table}[H]
\centering
\begin{tabular}{@{}lcc@{}}
\toprule
\textbf{Magnitud} & \textbf{Valor medido} & \textbf{Esperado} \\
\midrule
Período $T$              & \dato\ ms  & 1,000 ms \\
Ancho $t_{on}$           & \dato\ \us & --- \\
Carga de CPU $D$         & \dato\ \%  & --- \\
Ciclos equivalentes      & \dato      & --- \\
\bottomrule
\end{tabular}
\caption{M1: interrupción del SysTick.}
\end{table}

\textbf{Análisis.} El período medido fue \dato\ ms, lo que confirma que el
SysTick está configurado correctamente a 1 kHz según
\texttt{SYSTICK\_FREQ\_HZ}. La ISR consume \dato\ \us\ de cada milisegundo,
es decir un \dato\ \% de la CPU, dejando un \dato\ \% disponible para el
super-loop. El margen respecto del límite de 1 ms es de \dato\ \us, o sea que
la ISR podría crecer \dato\ veces antes de comprometer la base de tiempo.

% -----------------------------------------------------------------------------
\subsection{M2 --- Refresco del display}
% -----------------------------------------------------------------------------

\textbf{Configuración:} \texttt{TEST\_ISR\_TARGET = TEST\_ISR\_DISPLAY\_REFRESH}

\textbf{Objetivo.} Aislar el componente que se espera dominante dentro del
SysTick. \texttt{HAL\_Refresh\_Task()} hace \textit{bit-banging} de dos tramas
de 16 bits (una de \textit{blanking} anti-\textit{ghosting} y una útil) hacia
los shift-registers 74HC595: son del orden de un centenar de escrituras a GPIO
por cada tick.

\textbf{Procedimiento.} Medir el ancho en dos condiciones: con el display a
\textbf{brillo máximo} y a \textbf{brillo mínimo} (girando la perilla
presionada). El driver está escrito sin ramas por modo de operación, así que
se espera que el ancho \emph{no cambie}.

\captura{M2: pulso del refresco del display, a brillo máximo.}

\begin{table}[H]
\centering
\begin{tabular}{@{}lccc@{}}
\toprule
\textbf{Condición} & \textbf{Ancho} & \textbf{\% del SysTick (M1)} & \textbf{Ciclos} \\
\midrule
Brillo máximo (4 turnos)  & \dato[1.5cm] \us & \dato[1.2cm] \% & \dato[1.5cm] \\
Brillo mínimo (10 turnos) & \dato[1.5cm] \us & \dato[1.2cm] \% & \dato[1.5cm] \\
\bottomrule
\end{tabular}
\caption{M2: refresco del display en los dos extremos de brillo.}
\end{table}

\textbf{Análisis.} El refresco del display consume \dato\ \us, que representa
un \dato\ \% del tiempo total de la ISR del SysTick medido en M1: es,
efectivamente, el componente dominante.

La comparación entre brillo máximo y mínimo arroja una diferencia de
\dato\ \us. Este resultado es relevante porque el atenuador está implementado
por \textbf{reparto de turnos de multiplexado}: en los turnos ``vacíos'' se
manda la misma trama pero con los segmentos apagados. Es decir, el brillo no
se logra ejecutando más o menos código, sino cambiando \emph{qué} se envía.
La medición \dato{} (confirma / no confirma) esa hipótesis de diseño, lo que
implica que el costo de CPU del display es constante e independiente del
brillo elegido por el usuario.

% -----------------------------------------------------------------------------
\subsection{M3 --- FSM del encoder}
% -----------------------------------------------------------------------------

\textbf{Configuración:} \texttt{TEST\_ISR\_TARGET = TEST\_ISR\_ENCODER\_UPDATE}

\textbf{Objetivo.} Medir el segundo callback del SysTick.
\texttt{actualizarEncoder()} lee tres pines y avanza una máquina de estados de
antirrebote. Se espera un tiempo mucho menor al del display.

\textbf{Procedimiento.} Medir el ancho en reposo y, si se aprecia diferencia,
también mientras se gira la perilla (los caminos del \texttt{switch} que se
ejecutan son distintos).

\captura{M3: pulso de la actualización del encoder.}

\begin{table}[H]
\centering
\begin{tabular}{@{}lccc@{}}
\toprule
\textbf{Condición} & \textbf{Ancho} & \textbf{\% del SysTick (M1)} & \textbf{Ciclos} \\
\midrule
Perilla en reposo & \dato[1.5cm] \us & \dato[1.2cm] \% & \dato[1.5cm] \\
Girando la perilla & \dato[1.5cm] \us & \dato[1.2cm] \% & \dato[1.5cm] \\
\bottomrule
\end{tabular}
\caption{M3: FSM de antirrebote del encoder.}
\end{table}

\textbf{Análisis.} La FSM del encoder consume \dato\ \us\ por tick, un
\dato\ \% del presupuesto de la ISR. Es un costo despreciable frente al del
display, lo que justifica la decisión de sondear el encoder por
\textit{polling} desde el tick de 1 ms en lugar de usar interrupciones de pin:
se obtiene el antirrebote ``gratis'' con la base de tiempo que ya existe, sin
agregar fuentes de interrupción que compitan con el lector de tarjetas.

% -----------------------------------------------------------------------------
\subsection{M4 --- Base de tiempo en milisegundos}
% -----------------------------------------------------------------------------

\textbf{Configuración:} \texttt{TEST\_ISR\_TARGET = TEST\_ISR\_TIMER\_TICK}

\textbf{Objetivo.} Medir el tercer callback. \texttt{Timer\_Tick\_Callback()}
sólo incrementa un contador, así que debería ser el más corto de los tres.
Sirve de referencia del ``piso'' de costo de un callback.

\textbf{Procedimiento.} Puede requerir bajar mucho la base de tiempo del
osciloscopio. Si el pulso resulta demasiado angosto para medirlo con
confianza, dejarlo asentado como $<$ la resolución del equipo.

\captura{M4: pulso del contador de milisegundos.}

\begin{table}[H]
\centering
\begin{tabular}{@{}lcc@{}}
\toprule
\textbf{Magnitud} & \textbf{Valor medido} & \textbf{Ciclos} \\
\midrule
Ancho $t_{on}$ & \dato\ \us & \dato \\
\bottomrule
\end{tabular}
\caption{M4: callback de la base de tiempo.}
\end{table}

\textbf{Análisis.} El contador de milisegundos consume \dato\ \us. Al tratarse
de un incremento y una comparación, este valor representa aproximadamente el
\textbf{costo mínimo de un callback} en esta arquitectura, incluyendo la
llamada indirecta a través del puntero a función del \textit{scheduler}.

% -----------------------------------------------------------------------------
\subsection{M5 --- Balance del SysTick: el costo del despachador}
% -----------------------------------------------------------------------------

\textbf{Configuración:} no requiere medición nueva; se deriva de M1--M4.

\textbf{Objetivo.} Cuantificar cuánto cuesta el \emph{mecanismo} de callbacks.
La ISR del SysTick no ejecuta las tareas directamente: recorre un arreglo de
punteros a función y los invoca. Ese bucle es el precio de que los módulos
puedan registrarse sin que el MCAL los conozca.

\[
t_{\text{despachador}} \;=\; t_{\text{M1}} - \left( t_{\text{M2}} + t_{\text{M3}} + t_{\text{M4}} \right)
\]

\begin{table}[H]
\centering
\begin{tabular}{@{}lcc@{}}
\toprule
\textbf{Componente} & \textbf{Tiempo} & \textbf{\% del total} \\
\midrule
Refresco del display (M2) & \dato[1.5cm] \us & \dato[1.2cm] \% \\
FSM del encoder (M3)      & \dato[1.5cm] \us & \dato[1.2cm] \% \\
Base de tiempo (M4)       & \dato[1.5cm] \us & \dato[1.2cm] \% \\
\midrule
Suma de los callbacks     & \dato[1.5cm] \us & \dato[1.2cm] \% \\
Total medido (M1)         & \dato[1.5cm] \us & 100 \% \\
\midrule
\textbf{Despachador (diferencia)} & \dato[1.5cm] \us & \dato[1.2cm] \% \\
\bottomrule
\end{tabular}
\caption{M5: descomposición del tiempo de la ISR del SysTick.}
\end{table}

\textbf{Análisis.} El bucle despachador de \texttt{SysTick\_Handler()} agrega
\dato\ \us\ sobre la suma de las tareas, es decir un \dato\ \% de sobrecarga.
Este es el costo concreto de la abstracción: permite que \texttt{App\_Init()}
registre tareas con \texttt{HAL\_Scheduler\_AddTask()} sin que el MCAL sepa qué
hacen. El resultado \dato{} (justifica / no justifica) mantener el mecanismo
frente a la alternativa de llamar a las tres funciones en forma directa
dentro de la ISR.

% -----------------------------------------------------------------------------
\subsection{M6 --- Lector de tarjetas: captura de un bit}
% -----------------------------------------------------------------------------

\textbf{Configuración:} \texttt{TEST\_ISR\_TARGET = TEST\_ISR\_CARD\_CLK},
y luego \texttt{TEST\_ISR\_PORT\_DISPATCH}.

\textbf{Objetivo.} Caracterizar la ISR más crítica del sistema: se ejecuta una
vez por cada bit de la banda magnética y, si se pierde una sola ejecución, la
trama entera se corrompe. Además, comparando las dos configuraciones se
obtiene el costo del despachador de interrupciones de puerto
(\texttt{gpioIRQDispatch()}), que recorre los 32 pines buscando cuál generó la
interrupción.

\textbf{Procedimiento.} Pasar una tarjeta y capturar la ráfaga. Medir el ancho
de un pulso individual y el período entre pulsos consecutivos (que es la
velocidad real del \textit{swipe}). Repetir con las dos configuraciones.

\captura{M6: ráfaga de pulsos durante una pasada de tarjeta (uno por bit).}

\begin{table}[H]
\centering
\begin{tabular}{@{}lccc@{}}
\toprule
\textbf{Configuración} & \textbf{Ancho del pulso} & \textbf{Período entre bits} & \textbf{Ciclos} \\
\midrule
\texttt{CARD\_CLK} (sólo la ISR)      & \dato[1.5cm] \us & \dato[1.5cm] \us & \dato[1.5cm] \\
\texttt{PORT\_DISPATCH} (con despacho) & \dato[1.5cm] \us & \dato[1.5cm] \us & \dato[1.5cm] \\
\midrule
\textbf{Costo del despachador}         & \dato[1.5cm] \us & --- & \dato[1.5cm] \\
\bottomrule
\end{tabular}
\caption{M6: ISR de captura de bit, con y sin el despachador de puerto.}
\end{table}

\textbf{Análisis.} El cuerpo de \texttt{IRQ\_CardClk()} --- una lectura de pin,
un almacenamiento en el arreglo y una comparación --- consume \dato\ \us. Al
medir la interrupción completa con el despachador incluido, el tiempo sube a
\dato\ \us: el mecanismo genérico de \texttt{gpio.c} agrega \dato\ \us, un
\dato\ \% adicional.

Este resultado es interesante porque el despachador recorre 32 posiciones
aunque sólo una tenga un \textit{flag} activo. \dato{} (Se justifica / No se
justifica) frente a la alternativa de que cada driver escriba su propio
\texttt{PORTB\_IRQHandler}, considerando que a cambio se obtiene que cualquier
módulo pueda registrar un callback por pin sin tocar el MCAL.

El período entre bits medido (\dato\ \us) corresponde a la velocidad con la
que se pasó la tarjeta. La relación entre el tiempo de ISR y ese período
determina el margen del sistema, que se analiza en M8.

% -----------------------------------------------------------------------------
\subsection{M7 --- Lector de tarjetas: apertura y cierre del \textit{swipe}}
% -----------------------------------------------------------------------------

\textbf{Configuración:} \texttt{TEST\_ISR\_TARGET = TEST\_ISR\_CARD\_ENABLE}

\textbf{Objetivo.} Medir la ISR del pin ENABLE, que se ejecuta dos veces por
pasada (flanco de bajada al empezar, de subida al terminar). Se espera que sea
notablemente más larga que la de CLOCK, porque llama a \texttt{gpioIRQ()} para
armar y desarmar la interrupción de CLOCK, lo que implica configurar registros
(PCR, ISFR, NVIC).

\textbf{Procedimiento.} Pasar una tarjeta y capturar los dos pulsos. Si es
posible, provocar un \textbf{rebote} apoyando la tarjeta suavemente sin
completar la pasada, para observar el caso que descarta el
\textit{driver} por no alcanzar \texttt{MIN\_VALID\_BITS}.

\captura{M7: pulsos de apertura y cierre de la ventana de captura.}

\begin{table}[H]
\centering
\begin{tabular}{@{}lcc@{}}
\toprule
\textbf{Evento} & \textbf{Ancho} & \textbf{Ciclos} \\
\midrule
Flanco de bajada (apertura)  & \dato[1.5cm] \us & \dato[1.5cm] \\
Flanco de subida (cierre)    & \dato[1.5cm] \us & \dato[1.5cm] \\
Comparación con \texttt{CARD\_CLK} (M6) & \dato[1.5cm]$\times$ & --- \\
\bottomrule
\end{tabular}
\caption{M7: ISR de ENABLE del lector.}
\end{table}

\textbf{Análisis.} La ISR de ENABLE consume \dato\ \us, es decir \dato\ veces
el tiempo de la ISR de CLOCK medida en M6. La diferencia se explica porque
\texttt{IRQ\_CardEnable()} reconfigura la interrupción de CLOCK mediante
\texttt{gpioIRQ()}, operación que implica lectura-modificación-escritura sobre
el registro PCR del puerto, limpieza del \textit{flag} en ISFR y habilitación
en el NVIC.

Esta asimetría no representa un problema: la ISR larga ocurre sólo dos veces
por pasada y en momentos en que no hay bits llegando, mientras que la ISR que
se repite decenas de veces es la corta. \textbf{El diseño pone el trabajo
pesado donde no hay presión de tiempo}, que es exactamente el criterio
correcto.

% -----------------------------------------------------------------------------
\subsection{M8 --- Latencia de interrupción y validación de prioridades}
% -----------------------------------------------------------------------------

\textbf{Configuración:} \texttt{TEST\_ISR\_TARGET = TEST\_ISR\_CARD\_CLK},
\textbf{dos canales}.

\begin{itemize}
  \item CH1: PTB10 (CLOCK del lector) --- el flanco físico.
  \item CH2: PTC17 (pin de medición) --- la respuesta del software.
\end{itemize}

\textbf{Objetivo.} Éste es el ensayo más importante del conjunto. Mide el
tiempo que transcurre entre el flanco de bajada de CLOCK y el momento en que
el CPU efectivamente entra a la ISR. Ese retardo es la \textbf{latencia de
interrupción}, y su valor máximo determina si el sistema puede perder bits.

Además, permite validar experimentalmente una decisión de diseño concreta: en
\texttt{App\_Init()} se le baja la prioridad al SysTick
(\texttt{APP\_NVIC\_PRIO\_SYSTICK = 1}) para que la interrupción de PORTB
(prioridad 0) pueda interrumpirlo. Sin esa medida, un flanco de CLOCK que
llegara mientras corre la ISR del SysTick tendría que esperar a que ésta
termine.

\textbf{Procedimiento.} Disparar por CH1 (flanco de bajada). Medir el retardo
CH1$\rightarrow$CH2 con la función de medición de retardo entre canales.
Capturar \textbf{muchas} pasadas y quedarse con el valor mínimo y el máximo
observados: lo que importa es el peor caso.

\captura{M8: flanco de CLOCK (CH1) contra entrada a la ISR (CH2).}

\begin{table}[H]
\centering
\begin{tabular}{@{}lcc@{}}
\toprule
\textbf{Magnitud} & \textbf{Valor} & \textbf{Ciclos} \\
\midrule
Latencia mínima observada & \dato[1.5cm] \us & \dato[1.5cm] \\
Latencia máxima observada & \dato[1.5cm] \us & \dato[1.5cm] \\
Duración de la ISR del SysTick (M1) & \dato[1.5cm] \us & \dato[1.5cm] \\
\bottomrule
\end{tabular}
\caption{M8: latencia de interrupción del lector.}
\end{table}

\textbf{Análisis.} La latencia observada se mantuvo entre \dato\ \us\ y
\dato\ \us. El dato decisivo es la comparación entre la \textbf{latencia
máxima} y la \textbf{duración de la ISR del SysTick} medida en M1
(\dato\ \us):

\begin{itemize}
  \item Si la latencia máxima es del orden de unos pocos microsegundos y
        \emph{claramente menor} que la ISR del SysTick, queda demostrado que
        PORTB efectivamente \textbf{interrumpe} al SysTick, y la asignación de
        prioridades funciona como se diseñó.
  \item Si en cambio la latencia máxima se acercara a la duración de la ISR
        del SysTick, significaría que los flancos de CLOCK están esperando a
        que el SysTick termine, y el sistema estaría al borde de perder bits
        en las pasadas rápidas.
\end{itemize}

El resultado obtenido indica que \dato[6cm], lo que \dato{} (valida / invalida)
la elección de \texttt{APP\_NVIC\_PRIO\_SYSTICK = 1}.

\textbf{Ensayo complementario recomendado.} Repetir la medición forzando el
peor caso: pasar la tarjeta repetidamente hasta capturar una pasada en la que
un flanco de CLOCK coincida con la ejecución del SysTick. Si las prioridades
están bien asignadas, no debería observarse ningún ensanchamiento de la
latencia.

% -----------------------------------------------------------------------------
\subsection{M9 --- Super-loop}
% -----------------------------------------------------------------------------

\textbf{Configuración:} \texttt{TEST\_ISR\_TARGET = TEST\_ISR\_APP\_LOOP}

\textbf{Objetivo.} \texttt{App\_Run()} no es una interrupción, pero su período
determina la capacidad de respuesta de la aplicación. Además permite verificar
un contrato explícito documentado en \texttt{encoder.h}: el evento del encoder
vive como máximo 1 ms, por lo que \emph{el super-loop debe dar una vuelta
completa en menos de 1 ms} o se pierden pulsaciones.

\textbf{Procedimiento.} Medir período y ancho en reposo. Repetir durante una
pasada de tarjeta, cuando el lazo además ejecuta la decodificación en
\texttt{getCardData()}.

\captura{M9: vueltas del super-loop.}

\begin{table}[H]
\centering
\begin{tabular}{@{}lccc@{}}
\toprule
\textbf{Condición} & \textbf{Período} & \textbf{Ancho} & \textbf{Vueltas/s} \\
\midrule
Reposo (StandBy)          & \dato[1.5cm] \us & \dato[1.5cm] \us & \dato[1.5cm] \\
Decodificando una tarjeta & \dato[1.5cm] \us & \dato[1.5cm] \us & \dato[1.5cm] \\
\bottomrule
\end{tabular}
\caption{M9: período del lazo principal.}
\end{table}

\textbf{Análisis.} En reposo el super-loop completa una vuelta cada
\dato\ \us, es decir \dato\ vueltas por segundo. Como ese período es
\dato{} (menor / mayor) a 1 ms, el contrato de
\texttt{getEncoderEvent()} \dato{} (se cumple / no se cumple) y no se pierden
eventos del encoder.

Durante la decodificación de una tarjeta el período sube a \dato\ \us, porque
el lazo ejecuta además \texttt{getCardData()}, que recorre los bits crudos
buscando el \textit{Start Sentinel} y verifica la paridad de cada carácter.
Aun en ese caso el período se mantiene \dato{} (por debajo / por encima) del
límite de 1 ms. Esto confirma la decisión de diseño de \textbf{separar la
captura de la decodificación}: la ISR sólo guarda bits, y el trabajo costoso
se hace en el lazo principal, donde un retardo no causa pérdida de datos.

\newpage
% =============================================================================
\section{Análisis integrador}
% =============================================================================

\subsection{Resumen de tiempos}

\begin{table}[H]
\centering
\begin{tabular}{@{}llccc@{}}
\toprule
\textbf{Ensayo} & \textbf{Rutina} & \textbf{Ancho} & \textbf{Período} & \textbf{Duty} \\
\midrule
M1 & \texttt{SysTick\_Handler}   & \dato[1.3cm] \us & \dato[1.3cm] & \dato[1cm] \% \\
M2 & \texttt{HAL\_Refresh\_Task} & \dato[1.3cm] \us & 1 ms         & \dato[1cm] \% \\
M3 & \texttt{actualizarEncoder}  & \dato[1.3cm] \us & 1 ms         & \dato[1cm] \% \\
M4 & \texttt{Timer\_Tick\_Callback} & \dato[1.3cm] \us & 1 ms      & \dato[1cm] \% \\
M6 & \texttt{IRQ\_CardClk}       & \dato[1.3cm] \us & \dato[1.3cm] & --- \\
M6 & \texttt{gpioIRQDispatch}    & \dato[1.3cm] \us & \dato[1.3cm] & --- \\
M7 & \texttt{IRQ\_CardEnable}    & \dato[1.3cm] \us & 2 por pasada & --- \\
M9 & \texttt{App\_Run}           & \dato[1.3cm] \us & \dato[1.3cm] & \dato[1cm] \% \\
\bottomrule
\end{tabular}
\caption{Resumen de todos los tiempos medidos.}
\end{table}

\subsection{Presupuesto de CPU}

Con el sistema en reposo, el reparto del tiempo de procesador resulta:

\begin{table}[H]
\centering
\begin{tabular}{@{}lc@{}}
\toprule
\textbf{Consumidor} & \textbf{\% de CPU} \\
\midrule
ISR del SysTick (display + encoder + timer + despacho) & \dato[1.5cm] \% \\
Lector de tarjetas (sólo durante una pasada)           & \dato[1.5cm] \% \\
Super-loop (aplicación)                                 & \dato[1.5cm] \% \\
\midrule
\textbf{CPU disponible}                                 & \dato[1.5cm] \% \\
\bottomrule
\end{tabular}
\caption{Distribución de la carga de procesador.}
\end{table}

\subsection{Costo de las abstracciones de la arquitectura}

Dos de los ensayos cuantifican directamente el precio de las decisiones de
diseño en capas:

\begin{table}[H]
\centering
\begin{tabular}{@{}lcc@{}}
\toprule
\textbf{Abstracción} & \textbf{Sobrecarga} & \textbf{Qué se obtiene a cambio} \\
\midrule
\begin{tabular}[t]{@{}l@{}}Despachador de callbacks\\del SysTick (M5)\end{tabular}
  & \dato[1.3cm] \us
  & \begin{tabular}[t]{@{}l@{}}Los módulos se registran solos;\\el MCAL no los conoce.\end{tabular} \\
\addlinespace
\begin{tabular}[t]{@{}l@{}}Despachador de IRQ\\de puerto (M6)\end{tabular}
  & \dato[1.3cm] \us
  & \begin{tabular}[t]{@{}l@{}}Callback por pin sin que cada\\driver escriba su handler.\end{tabular} \\
\bottomrule
\end{tabular}
\caption{Costo medido de los mecanismos genéricos.}
\end{table}

\subsection{Conclusiones}

\begin{enumerate}
  \item \textbf{Restricción de tiempo real.} La ISR del SysTick consume
        \dato\ \% del período de 1 ms, con un margen de \dato\ \us. El sistema
        \dato{} (cumple / no cumple) su restricción fundamental.

  \item \textbf{Componente dominante.} El refresco del display representa el
        \dato\ \% del tiempo de interrupción. Cualquier optimización futura
        debería atacar el \textit{bit-banging} hacia los shift-registers ---
        por ejemplo, migrándolo a SPI por hardware ---, ya que el resto de las
        tareas es despreciable en comparación.

  \item \textbf{Asignación de prioridades.} La latencia máxima del lector
        (\dato\ \us) frente a la duración de la ISR del SysTick (\dato\ \us)
        \dato{} (confirma / refuta) que bajar la prioridad del SysTick fue la
        decisión correcta.

  \item \textbf{Separación captura/decodificación.} La ISR por bit consume
        \dato\ \us\ contra un período entre bits de \dato\ \us, lo que da un
        margen de \dato$\times$. La tarjeta podría pasarse \dato\ veces más
        rápido antes de comprometer la captura.

  \item \textbf{Costo de la arquitectura.} Los mecanismos genéricos suman
        \dato\ \us\ de sobrecarga total, un \dato\ \% del presupuesto. El
        equipo considera que este costo \dato{} (se justifica / no se
        justifica) frente a la ganancia en modularidad y portabilidad.

  \item \textbf{Efecto de la optimización.} Todas las mediciones corresponden
        a \texttt{-O0}. Se estima que con \texttt{-O2} los tiempos se
        reducirían \dato{}, aunque las proporciones relativas se mantendrían.
\end{enumerate}

\appendix
\newpage
% =============================================================================
\section{Opciones de instrumentación disponibles}
% =============================================================================

\begin{table}[H]
\centering
\begin{tabular}{@{}llc@{}}
\toprule
\textbf{Valor de \texttt{TEST\_ISR\_TARGET}} & \textbf{Rutina instrumentada} & \textbf{Ensayo} \\
\midrule
\texttt{TEST\_ISR\_NONE}            & --- (instrumentación apagada)   & ---  \\
\texttt{TEST\_ISR\_SYSTICK\_TOTAL}  & \texttt{SysTick\_Handler}       & M1   \\
\texttt{TEST\_ISR\_DISPLAY\_REFRESH}& \texttt{HAL\_Refresh\_Task}     & M2   \\
\texttt{TEST\_ISR\_ENCODER\_UPDATE} & \texttt{actualizarEncoder}      & M3   \\
\texttt{TEST\_ISR\_TIMER\_TICK}     & \texttt{Timer\_Tick\_Callback}  & M4   \\
\texttt{TEST\_ISR\_PORT\_DISPATCH}  & \texttt{gpioIRQDispatch}        & M6   \\
\texttt{TEST\_ISR\_CARD\_CLK}       & \texttt{IRQ\_CardClk}           & M6, M8 \\
\texttt{TEST\_ISR\_CARD\_ENABLE}    & \texttt{IRQ\_CardEnable}        & M7   \\
\texttt{TEST\_ISR\_APP\_LOOP}       & \texttt{App\_Run}               & M9   \\
\bottomrule
\end{tabular}
\caption{Opciones definidas en \texttt{source/test\_isr\_cfg.h}.}
\end{table}

\section{Pines de interés para el osciloscopio}

\begin{table}[H]
\centering
\begin{tabular}{@{}lll@{}}
\toprule
\textbf{Señal} & \textbf{Pin} & \textbf{Uso en las mediciones} \\
\midrule
Pin de instrumentación & PTC17 & Todas (CH2 en M8) \\
CLOCK del lector       & PTB10 & CH1 en M8 (flanco de referencia) \\
ENABLE del lector      & PTB9  & Enmarcar la pasada completa \\
DATA del lector        & PTB11 & Verificar el \textit{setup} del dato \\
LATCH de los 74HC595   & PTC4  & Contar tramas (2 por turno: blanking + útil) \\
\bottomrule
\end{tabular}
\caption{Puntos de medición disponibles en la placa.}
\end{table}

\end{document}
